\documentclass{report}

\input{~/dev/latex/template/preamble.tex}
\input{~/dev/latex/template/macros.tex}

\title{\Huge{}}
\author{\huge{Nathan Warner}}
\date{\huge{}}
\fancyhf{}
\rhead{}
\fancyhead[R]{\itshape Warner} % Left header: Section name
\fancyhead[L]{\itshape\leftmark}  % Right header: Page number
\cfoot{\thepage}
\renewcommand{\headrulewidth}{0pt} % Optional: Removes the header line
%\pagestyle{fancy}
%\fancyhf{}
%\lhead{Warner \thepage}
%\rhead{}
% \lhead{\leftmark}
%\cfoot{\thepage}
%\setborder
% \usepackage[default]{sourcecodepro}
% \usepackage[T1]{fontenc}

% Change the title
\hypersetup{
    pdftitle={Math 4}
}

\begin{document}
    % \maketitle
        \begin{titlepage}
       \begin{center}
           \vspace*{1cm}
    
           \textbf{Math 4: Numerical Linear Algebra}
    
           \vspace{0.5cm}
            
                
           \vspace{1.5cm}
    
           \textbf{Nathan Warner}
    
           \vfill
                
                
           \vspace{0.8cm}
         
           \includegraphics[width=0.4\textwidth]{~/niu/seal.png}
                
           Computer Science \\
           Northern Illinois University\\
           United States\\
           
                
       \end{center}
    \end{titlepage}
    \tableofcontents
    \pagebreak 
    \unsect{Numerical Linear Algebra}
    \bigbreak \noindent 
    \subsection{Introduction}
    \bigbreak \noindent 
    \begin{itemize}
        \item \textbf{Matrix Notation}: For a matrix $A \in \mathbb{R}^{m\times n} $, we say
            \begin{align*}
                A = (a_{ij}) = \begin{bmatrix}
                    a_{11} & a_{12} & \cdots & a_{1n} \\
                    a_{21} & a_{22} & \cdots & a_{2n} \\
                    \vdots & \vdots & \ddots & \vdots \\
                    a_{m1} & a_{m2} & \cdots & a_{mn}
                \end{bmatrix}
            \end{align*}
            with $a_{ij} \in \mathbb{R} $.
        \item \textbf{Vector notation}: For a vector $x \in \mathbb{R}^{n} $ (or $\mathbb{R}^{n\times 1} $), we have
            \begin{align*}
                x = \begin{bmatrix}
                    x_{1} \\ x_{2} \\ \vdots \\ x_{n}
                \end{bmatrix}
            \end{align*}
            for $x_{i} \in \mathbb{R} $.

        \item \textbf{Submatrix notation (rows)}: 
            \[
                A(i,:) \in \mathbb{R}^{1 \times n} \;\;\Longleftrightarrow\;\;
                A(i,:) = \bigl[\, a_{i1} \; a_{i2} \; \cdots \; a_{in} \,\bigr].
            \]
        \item \textbf{Submatrix notation (columns)}: 
            \[
                A(:,j) \in \mathbb{R}^{m \times 1} \;\;\Longleftrightarrow\;\;
                A(:,j) =
                \begin{bmatrix}
                    a_{1j} \\
                    a_{2j} \\
                    \vdots \\
                    a_{mj}
                \end{bmatrix}.
            \]
        \item \textbf{Sparse Matrix}: A sparse matrix or sparse array is a matrix in which most of the elements are zero. There is no strict definition regarding the proportion of zero-value elements for a matrix to qualify as sparse but a common criterion is that the number of non-zero elements is roughly equal to the number of rows or columns.
        \item \textbf{Dense Matrix}: if most of the elements are non-zero, the matrix is considered dense
        \item \textbf{Sparsity}: The number of zero-valued elements divided by the total number of elements is sometimes referred to as the sparsity of the matrix.
        \item \textbf{Band Matrix}: a band matrix or banded matrix is a sparse matrix whose non-zero entries are confined to a diagonal band, comprising the main diagonal and zero or more diagonals on either side.
            \bigbreak \noindent 
            \[
                A(i_{1}:i_{2},:) \in \mathbb{R}^{(i_{2}-i_{1}+1) \times n}
                \;\;\Longleftrightarrow\;\;
                A(i_{1}:i_{2},:) =
                \begin{bmatrix}
                    a_{i_{1}1} & a_{i_{1}2} & \cdots & a_{i_{1}n} \\
                    a_{i_{1}+1,1} & a_{i_{1}+1,2} & \cdots & a_{i_{1}+1,n} \\
                    \vdots & \vdots & \ddots & \vdots \\
                    a_{i_{2}1} & a_{i_{2}2} & \cdots & a_{i_{2}n}
                \end{bmatrix}.
            \]

            \[
                A(:,j_{1}:j_{2}) \in \mathbb{R}^{m \times (j_{2}-j_{1}+1)}
                \;\;\Longleftrightarrow\;\;
                A(:,j_{1}:j_{2}) =
                \begin{bmatrix}
                    a_{1j_{1}} & a_{1,j_{1}+1} & \cdots & a_{1j_{2}} \\
                    a_{2j_{1}} & a_{2,j_{1}+1} & \cdots & a_{2j_{2}} \\
                    \vdots & \vdots & \ddots & \vdots \\
                    a_{mj_{1}} & a_{m,j_{1}+1} & \cdots & a_{mj_{2}}
                \end{bmatrix}.
            \]
            Where
            \[
                \begin{aligned}
                    A(i_{1}:i_{2},:) &:\;\; \text{all rows between } i_{1} \text{ and } i_{2}, \;\text{across all columns}, \\[6pt]
                    A(:,j_{1}:j_{2}) &:\;\; \text{all columns between } j_{1} \text{ and } j_{2}, \;\text{across all rows}.
                \end{aligned}
            \]
        \item \textbf{Transposition}: $\mathbb{R}^{m\times n} \to \mathbb{R}^{m\times n} $
            \begin{align*}
                C = A^{\top} \iff c_{ij} = a_{ji}
            .\end{align*}
        \item \textbf{Addition} $\quad (\mathbb{R}^{m \times n} \times \mathbb{R}^{m \times n} \;\to\; \mathbb{R}^{m \times n})$
            \[
                C = A + B 
                \;\;\Longrightarrow\;\; 
                c_{ij} = a_{ij} + b_{ij}.
            \]

        \item \textbf{Scalar-matrix Multiplication: } $\quad (\mathbb{R} \times \mathbb{R}^{m \times n} \;\to\; \mathbb{R}^{m \times n})$
            \[
                C = \alpha A 
                \;\;\Longrightarrow\;\; 
                c_{ij} = \alpha a_{ij}.
            \]

        \item \textbf{Matrix-matrix Multiplication: } $\quad (\mathbb{R}^{m \times p} \times \mathbb{R}^{p \times n} \;\to\; \mathbb{R}^{m \times n})$
            \[
                C = AB 
                \;\;\Longrightarrow\;\; 
                c_{ij} = \sum_{k=1}^{p} a_{ik} b_{kj}.
            \]
        \item \textbf{Matrix-vector Multiplication: } $\quad (\mathbb{R}^{m \times n} \times \mathbb{R}^n \;\to\; \mathbb{R}^m)$
            \[
                y = Ax 
                \;\;\Longrightarrow\;\; 
                y_i = \sum_{j=1}^{n} a_{ij} x_j.
            \]

        \item \textbf{Inner product (or dot product): } $\quad (\mathbb{R}^n \times \mathbb{R}^n \;\to\; \mathbb{R})$
            \[
                c = x^T y 
                \;\;\Longrightarrow\;\; 
                c = \sum_{i=1}^{n} x_i y_i.
            \]

            \item \textbf{Outer product: } $\quad (\mathbb{R}^m \times \mathbb{R}^n \;\to\; \mathbb{R}^{m \times n})$
            \[
                C = x y^T 
                \;\;\Longrightarrow\;\; 
                c_{ij} = x_i y_j.
            \]
        \item \textbf{Flops}: A flop is a floating-point operation between numbers stored in a floating-point format on a computer.
            \bigbreak \noindent 
            If $x$ and $y$ are numbers stored in a floating point format, then the following operations are each one flop
            \begin{align*}
                x + y \quad x - y \quad xy \quad x / y
            .\end{align*}







    \end{itemize}



    \pagebreak \bigbreak \noindent 
    \subsection{Gaussian Elimination and its variants}
    \bigbreak \noindent 
    \subsubsection{Matrix Multiplication}
    \bigbreak \noindent 
    \begin{itemize}
        \item \textbf{Matrix Multiplication}:        
            In general, if $A$ is a real matrix with $m$ rows and $n$ columns, and $x$ is a real vector with $n$ entries, then
            \[
                A = 
                \begin{bmatrix}
                    a_{11} & \cdots & a_{1n} \\
                    \vdots & \ddots & \vdots \\
                    a_{m1} & \cdots & a_{mn}
                \end{bmatrix}
                \in \mathbb{R}^{m \times n}
                \quad \text{and} \quad
                x =
                \begin{bmatrix}
                    x_{1} \\
                    \vdots \\
                    x_{n}
                \end{bmatrix}
                \in \mathbb{R}^n.
            \]
            If $b = Ax$, then $b \in \mathbb{R}^m$ and
            \[
                b_i = \sum_{j=1}^n a_{ij} x_j
                = a_{i1}x_1 + \cdots + a_{in}x_n, 
                \quad i = 1, \ldots, m.
            \]
            Thus, $b_i$ is the \textbf{inner-product} between the $i$-row of $A$, 
            \[
                A(i,:) = [a_{i1} \;\; \cdots \;\; a_{in}], \quad (i = 1, \ldots, m)
            \]
            and the vector $x$.
            \bigbreak \noindent 
            Also,
            \[
                b = A(:,1) x_1 + \cdots + A(:,n)x_n,
            \]
            so $b$ is a \textbf{linear combination} of the columns of $A$, i.e.,
            \[
                A(:,j) = 
                \begin{bmatrix}
                    a_{1j} \\
                    a_{2j} \\
                    \vdots \\
                    a_{mj}
                \end{bmatrix},
                \quad j = 1, \ldots, n.
            \]
        \item \textbf{Matrix-Matrix Multiplication}:
            Let $A \in \mathbb{R}^{m \times n}$ and $X \in \mathbb{R}^{n \times p}$.
            \bigbreak \noindent 
            If $B = AX$ then $B \in \mathbb{R}^{m \times p}$ and
            \[
                b_{ij} = \sum_{k=1}^n a_{ik} x_{kj}, \quad i = 1, \ldots, m, \;\; j = 1, \ldots, p.
            \]
            That is, $b_{ij}$ is the inner-product between row $i$ of $A$ and column $j$ of $X$.
            \bigbreak \noindent 
            Also, each column of $B$ is a linear combination of the columns of $A$.
            \bigbreak \noindent 
            Total flops required for matrix multiplication is
            \[
                \sum_{i=1}^m \sum_{j=1}^p \sum_{k=1}^n 2 = 2mnp.
            \]
            If $A, X \in \mathbb{R}^{n \times n}$, then computing $B = AX$ requires $2n^3 = O(n^3)$ flops.
        \item \textbf{Block Matrices}:
            Partition $A \in \mathbb{R}^{m \times n}$ and $X \in \mathbb{R}^{n \times p}$ into blocks:
            \[
                A =
                \begin{blockarray}{ccc}
   & n_1 & n_2 \\
   \begin{block}{c[cc]}
       m_1 & A_{11} & A_{12} \\
       m_2 & A_{21} & A_{22} \\
   \end{block}
                \end{blockarray}
                \quad,\qquad
                X =
                \begin{blockarray}{ccc}
   & p_1 & p_2 \\
   \begin{block}{c[cc]}
       n_1 & X_{11} & X_{12} \\
       n_2 & X_{21} & X_{22} \\
   \end{block}
                \end{blockarray}
            \]
                    where $n = n_1 + n_2$, $m = m_1 + m_2$, and $p = p_1 + p_2$.
                    \bigbreak \noindent 
                    If $B = AX$, and 
                    \begin{align*}
                        B = 
                        \begin{blockarray}{ccc}
                            & p_{1} & p_{2} \\
                            \begin{block}{c[cc]}
                                m_{1} & B_{11} & B_{12} \\
                                m_{2} & B_{21} & B_{22} \\
                            \end{block}
                        \end{blockarray}
                    ,\end{align*}
                    then
                    \[
                        \begin{blockarray}{cc}
  &  \\
  \begin{block}{[cc]}
      B_{11} & B_{12} \\
      B_{21} & B_{22} \\
  \end{block}
                        \end{blockarray}
                        =
                        B = AX =
                        \begin{blockarray}{cc}
  &  \\
  \begin{block}{[cc]}
      A_{11} & A_{12} \\
      A_{21} & A_{22} \\
  \end{block}
                        \end{blockarray}
                        \begin{blockarray}{cc}
  &  \\
  \begin{block}{[cc]}
      X_{11} & X_{12} \\
      X_{21} & X_{22} \\
  \end{block}
                        \end{blockarray}
                    \]

                    \[
                        =
                        \begin{blockarray}{cc}
  &  \\
  \begin{block}{[cc]}
      A_{11}X_{11} + A_{12}X_{21} & A_{11}X_{12} + A_{12}X_{22} \\
      A_{21}X_{11} + A_{22}X_{21} & A_{21}X_{12} + A_{22}X_{22} \\
  \end{block}
                        \end{blockarray}
                    \]

                    That is,
                    \[
                        B_{ij} = \sum_{k=1}^2 A_{ik} X_{kj}, 
                        \qquad i,j = 1,2.
                    \]




    \end{itemize}

    \pagebreak 
    \subsubsection{Systems of Linear Equations}
    \begin{itemize}
        \item \textbf{Systems of linear equations}: Let $A \in \mathbb{R}^{n\times n},\ b \in \mathbb{R}^{m}$, our goal is to find $x \in \mathbb{R}^{m}$ such that $Ax = b$
        \item \textbf{Singularity}: A \textbf{singular matrix} is a square matrix that does not have an inverse.
            \bigbreak \noindent 
            A \textbf{nonsingular} matrix is a square matrix that does have an inverse.
            \bigbreak \noindent 
            The following are equivalent
            \begin{itemize}
                \item $Ax = b$ has a unique solution
                \item $\det(A)\ne 0$
                \item $A^{-1}$ exists
                \item There is no nonzero vector $y \in \mathbb{R}^{m}$ such that $Ay=0 $
            \end{itemize}
            If any one of the following are true, they all are true, and $A$ is non-singular
        \item \textbf{Solution to $Ax = b$}: If $A$ is nonsingular, then $A^{-1}$ exists, and
            \begin{align*}
                x = A^{-1}b
            .\end{align*}
            Which is the unique solution to $Ax=b$
            \bigbreak \noindent 
            \textbf{Note:} Practically, it is not wise to compute $A^{-1}$, as this can be expensive.
    \end{itemize}

    \pagebreak 
    \subsubsection{Triangular systems}
    \begin{itemize}
        \item \textbf{Upper triangular}: A square matrix $A = U$ of the form 
            \begin{align*}
                A = U = \begin{bmatrix}
                    u_{11} & u_{12} & \cdots & u_{1n} \\
                    0 & u_{22} & \cdots & u_{2n} \\
                    \vdots & \vdots & \ddots & \vdots \\
                    0 & 0 & \cdots & u_{nn}
                \end{bmatrix}
            \end{align*}
            is called \textbf{upper triangular}
        \item \textbf{Lower triangular}: A square matrix $A = L$
            \begin{align*}
                A = L = \begin{bmatrix}
                    \ell_{11} & 0 & \cdots & 0 \\
                    \ell_{21} & \ell_{22} & \cdots & 0 \\
                    \vdots & \vdots & \ddots & \vdots \\
                    \ell_{n1} & \ell_{n2} & \cdots & \ell_{nn}
                \end{bmatrix}
            \end{align*}
            is called \textbf{lower triangular}
        \item \textbf{Solutions to triangular systems}: Consider the system
            \begin{align*}
               \begin{bmatrix}
                    \ell_{11} & 0 & \cdots & 0 \\
                    \ell_{21} & \ell_{22} & \cdots & 0 \\
                    \vdots & \vdots & \ddots & \vdots \\
                    \ell_{n1} & \ell_{n2} & \cdots & \ell_{nn}
                \end{bmatrix} \begin{bmatrix}
                    x_{1} \\ x_{2} \\ \vdots \\ x_{n}
                \end{bmatrix}
                =
                \begin{bmatrix}
                    b_{1} \\ b_{2} \\ \vdots \\ b_{n}
                \end{bmatrix}
            .\end{align*}
            So,
            \begin{align*}
                \ell_{11}x_{1} &= b_{1} \\
                \ell_{21}x_{1} + \ell_{22}x_{2} &= b_{2} \\
                \vdots \\
                \ell_{n1}x_{1} + \ell_{n2}x_{2} + \cdots + \ell_{nn}x_{n} &= b_{n}
            .\end{align*}
            Then,
            \begin{align*}
                x_{1} &= \frac{b_{1}}{\ell_{11}}
            \end{align*}
            and,
            \begin{align*}
                \ell_{22}x_{2} &= b_{2} - \ell_{21}x_{1} \\
                \implies x_{2} &= \frac{b_{2}-\ell_{21}x_{1}}{\ell_{22}}
            .\end{align*}
            In general, we have
            \begin{align*}
                x_{i} = \frac{b_{i} - \sum_{j=1}^{i-1}\ell_{ij}x_{j}}{\ell_{ii}}
            \end{align*}
            for $i =1,2,...,n $. This method is called \textbf{Forward Substitution}.
            \bigbreak \noindent 
            A similar process is used on upper triangular matrices and is called \textbf{Backward Substitution}.
        \item \textbf{Counting flops for the forward substitution method}: We have 
                \begin{jlcode}
                for i = 1:n
                    for j=1:i-1
                        b[i] = b[i] - ell[i,j]b[j]
                    end
                    b[i] = b[i] / ell[i,i]
                end
                \end{jlcode}
            Thus, the count of flops is
            \begin{align*}
                n+\sum_{i=1}^{n}2(i-1) &= n+2 \sum_{i=1}^{n}(i-1) = n+2 \left( \sum_{i=1}^{n}i - \sum_{i=1}^{n} 1\right)  \\
                   &= n + 2 \left(\sum_{i=1}^{n}i - n\right) = n+ 2 \left(\frac{n(n+1)}{2}-n\right) \\
                   &=n + n^{2} - n = n^{2}
            \end{align*}
            So, forward substitution is $\mathcal{O}(n^{2})$
        \item \textbf{Column oriented forward substitution}: Suppose we have $Lx = b$ when $L$ is lower triangular, we split the matrix into the following blocks
            \begin{align*}
                \begin{bmatrix}
                    \ell_{11} & 0 \\
                    \hat{\ell} & \hat{L}
                \end{bmatrix}
                \begin{bmatrix}
                    x_{1} \\ \hat{x} 
                \end{bmatrix}
                = \begin{bmatrix}
                    b_{1} \\ \hat{b}
                \end{bmatrix}
            .\end{align*}
            With $\hat{\ell} \in \mathbb{R}^{n-1}$, $\hat{L} \in \mathbb{R}^{n-1 \times n-1} $, $\hat{x} \in \mathbb{R}^{n-1}$, $\ell_{11}, x_{1}, b_{1} \in \mathbb{R}$. Note that $\hat{L}$ is also lower triangular.
            \bigbreak \noindent 
            We have
            \begin{align*}
                \ell_{11}x_{1} &= b_{1} \implies x_{1} = \frac{b_{1}}{\ell_{11}} \\
                \hat{\ell}x_{1} + \hat{L}\hat{x} &= \hat{b} \implies \hat{L}\hat{x} = \hat{b} - \hat{\ell}x_{1}
            \end{align*}
            Thus, we reduced the dimension by one. We repeat this process for the remaining $x_{i}$. The process is 
            \begin{enumerate}
                \item Compute $x_{1} = \frac{b_{1}}{\ell_{11}} $
                \item Compute $\hat{b} - \hat{\ell}x_{1} = \tilde{b} \in \mathbb{R}^{n-1} $
                \item Find $\hat{L}x = \tilde{b} $
            \end{enumerate}
        \item \textbf{Counting flops for column oriented forward substitution}: Let $f_{n}$ be the flop count, we have
            \begin{align*}
                f_{n} = 1 + 2(n-1) + f_{n-1}
            .\end{align*}
            With
            \begin{align*}
                f_{n-1} = 1 + 2(n-2) + f_{n-2}
            .\end{align*}
            Until
            \begin{align*}
                f_{1} = 1 + 2 + f_{0}
            \end{align*}
            with $f_{0} = 0 $
            \bigbreak \noindent 
            So, 
            \begin{align*}
                f_{n} &= 1 + 2(n-1) + 1 + 2(n-2) + ... + 1 + 2(n-(n-1)) \\
                &= \sum_{i=1}^{n-1} 1 + 2(n-1) = n + 2n^{2} - 2\sum_{i=1}^{n}i \\
                &=...=n^{2}
            .\end{align*}
            Thus, column oriented forward substitution is also $\mathcal{O}(n^{2})$


            


            
    \end{itemize}

    \pagebreak 
    \unsect{Julia}
    \bigbreak \noindent 
    \subsection{Types}
    \bigbreak \noindent 
    \fig{.5}{./figures/tree.png}
    \begin{itemize}
        \item \textbf{Subtype constraint <:} $A$ <: $B$ means $A$ is a subtype of $B$
            \bigbreak \noindent 
            \begin{jlcode}
            Int <: Number #true
            \end{jlcode}
    \end{itemize}

    \pagebreak 
    \subsection{Functions}
    \begin{itemize}
        
    \end{itemize}

    \pagebreak 
    \subsection{Linear Algebra}
    \bigbreak \noindent 
    \subsubsection{Matrix creation and operations}
    \begin{itemize}
        \item \textbf{Array constructors}:
            \begin{itemize}
                \item Array\{T\}(undef, dims...)
                \item Vector\{T\}(undef, n)
                \item Matrix\{T\}(undef, m, n)
            \end{itemize}
        \item \textbf{Zeros/ones/fills}
            \begin{itemize}
                \item zeros(n), zeros(m,n)
                \item ones(n), ones(m,n)
                \item fill(x, dims...)
            \end{itemize}
        \item \textbf{Uniform ranges}:
            \begin{itemize}
                \item collect(1:n) $\to$ vector
                \item collect(1:m, 1:n) $\to$ matrix (grid)
            \end{itemize}
    \end{itemize}

    
\end{document}
